\documentclass[11pt,a4paper]{article}

% --- Encoding ---
\usepackage[utf8]{inputenc}
\usepackage[T1]{fontenc}
\usepackage[spanish,es-noshorthands]{babel}
\usepackage{lmodern}

% --- Layout ---
\usepackage[margin=2.5cm]{geometry}
\usepackage{parskip}
\setlength{\parindent}{0pt}
\setlength{\parskip}{0.5em plus 0.2em minus 0.1em}

% --- Tables and graphics ---
\usepackage{booktabs}
\usepackage{array}
\usepackage{enumitem}
\usepackage{xcolor}
\usepackage{graphicx}
\usepackage{tikz}
\usetikzlibrary{positioning,shapes.geometric,arrows.meta,fit,backgrounds,calc,decorations.pathreplacing}
\usepackage{caption}
\usepackage{subcaption}
\usepackage{amsmath,amssymb}

% --- Soporte para caracteres Unicode usados en formulas en texto plano ---
\usepackage{newunicodechar}
\newunicodechar{→}{\ensuremath{\rightarrow}}
\newunicodechar{≥}{\ensuremath{\geq}}
\newunicodechar{≤}{\ensuremath{\leq}}
\newunicodechar{×}{\ensuremath{\times}}
\newunicodechar{÷}{\ensuremath{\div}}
\newunicodechar{≈}{\ensuremath{\approx}}
\newunicodechar{⌈}{\ensuremath{\lceil}}
\newunicodechar{⌉}{\ensuremath{\rceil}}
\newunicodechar{✓}{\ensuremath{\checkmark}}

% --- Code listings ---
\usepackage{listings}
\definecolor{codegray}{rgb}{0.96,0.96,0.96}
\definecolor{codeborder}{rgb}{0.85,0.85,0.85}
\definecolor{codecomment}{rgb}{0.30,0.55,0.30}
\definecolor{codekeyword}{rgb}{0.10,0.30,0.70}
\definecolor{codestring}{rgb}{0.65,0.15,0.10}

\lstdefinestyle{cstyle}{
    language=C,
    backgroundcolor=\color{codegray},
    basicstyle=\ttfamily\footnotesize,
    breaklines=true,
    frame=single,
    rulecolor=\color{codeborder},
    showstringspaces=false,
    keywordstyle=\color{codekeyword}\bfseries,
    commentstyle=\color{codecomment}\itshape,
    stringstyle=\color{codestring},
    aboveskip=1em,
    belowskip=1em,
}

\lstdefinestyle{shellstyle}{
    language=bash,
    backgroundcolor=\color{codegray},
    basicstyle=\ttfamily\small,
    breaklines=true,
    frame=single,
    rulecolor=\color{codeborder},
    showstringspaces=false,
    keywordstyle=\color{codekeyword}\bfseries,
    commentstyle=\color{codecomment}\itshape,
    stringstyle=\color{codestring},
    aboveskip=1em,
    belowskip=1em,
    columns=fullflexible,
    upquote=true,
}

\lstdefinestyle{textstyle}{
    backgroundcolor=\color{codegray},
    basicstyle=\ttfamily\footnotesize,
    breaklines=true,
    frame=single,
    rulecolor=\color{codeborder},
    showstringspaces=false,
    aboveskip=1em,
    belowskip=1em,
    columns=fullflexible,
}

\lstset{style=textstyle}

% --- Definition / Idea / Important boxes ---
\usepackage[most]{tcolorbox}
\newtcolorbox{idea}[1][]{
    colback=orange!5!white,
    colframe=orange!60!black,
    fonttitle=\bfseries,
    title=Idea intuitiva,
    sharp corners=southwest,
    rounded corners=northwest,
    arc=2pt,
    boxrule=0.6pt,
    left=8pt,right=8pt,top=4pt,bottom=4pt,
    #1
}
\newtcolorbox{importante}[1][]{
    colback=red!5!white,
    colframe=red!50!black,
    fonttitle=\bfseries,
    title=Importante,
    sharp corners=southwest,
    rounded corners=northwest,
    arc=2pt,
    boxrule=0.6pt,
    left=8pt,right=8pt,top=4pt,bottom=4pt,
    #1
}
\newtcolorbox{respuesta}[1][]{
    colback=green!5!white,
    colframe=green!40!black,
    fonttitle=\bfseries,
    title=Respuesta,
    sharp corners=southwest,
    rounded corners=northwest,
    arc=2pt,
    boxrule=0.6pt,
    left=8pt,right=8pt,top=4pt,bottom=4pt,
    #1
}

% --- Hyperlinks ---
\usepackage[colorlinks=true,
            linkcolor=blue!60!black,
            urlcolor=blue!50!black,
            citecolor=blue!60!black]{hyperref}

% --- Color palette ---
\definecolor{c1}{HTML}{4C72B0}
\definecolor{c2}{HTML}{DD8452}
\definecolor{c3}{HTML}{55A868}
\definecolor{c4}{HTML}{8172B2}
\definecolor{cred}{HTML}{C44E52}
\definecolor{cfree}{HTML}{8C8C8C}

% --- Title block ---
\title{
    \vspace{-2.5em}
    \textbf{Recuperatorio del Parcial 1 --- 2025} \\[0.3em]
    \large Sistemas Operativos
}
\author{
    Tomás Rodeghiero \\
    \small UNRC --- Universidad Nacional de Río Cuarto
}
\date{2025}

\begin{document}

\maketitle

\begin{abstract}
\noindent
Resolución completa del Recuperatorio del Parcial 1 del año 2025 de
la cátedra de Sistemas Operativos (UNRC). El parcial consta de cinco
ejercicios: (i) descripción detallada del flujo de una
interrupción/excepción y el rol del \emph{trapframe}; (ii)
implementación en pseudo-código de un \texttt{lock} reentrante (que
evita \emph{self-lock} cuando el mismo dueño lo invoca repetidamente);
(iii) comparación entre los algoritmos de planificación FIFO y Round
Robin; (iv) verdadero/falso sobre paginación, segmentación,
\emph{copy-on-write} y \texttt{mmap}; (v) cálculos sobre una
arquitectura de 32 bits con páginas de 4\,KB, incluyendo construcción
de la tabla de páginas de un proceso y determinación de page faults.
Cada ejercicio se desarrolla con la estructura \textbf{Enunciado
$\to$ Análisis $\to$ Desarrollo $\to$ Respuesta}.
\end{abstract}

\tableofcontents
\bigskip

% =====================================================================
\section{Ejercicio 1 --- Interrupciones, excepciones y \emph{trapframe} (2 ptos)}
% =====================================================================

\textbf{Enunciado.} Describa de forma general los pasos que sigue la
CPU y el código de un SO en una interrupción o excepción hasta su
retorno. En particular, describa el rol del \emph{trapframe}.

\subsection*{Análisis: qué es un \emph{trap}}

Una \textbf{trap} es cualquier evento que interrumpe la ejecución
normal de un proceso y transfiere el control al kernel. Hay tres
familias:

\begin{itemize}[leftmargin=*,nosep]
    \item \textbf{Interrupción}: viene de un dispositivo externo
        (timer, disco, red, teclado). Es \emph{asincrónica}: ocurre
        en cualquier punto de la ejecución.
    \item \textbf{Excepción}: se origina en la propia instrucción
        que está ejecutando la CPU (page fault, división por cero,
        instrucción ilegal, acceso a memoria fuera de rango).
        \emph{Sincrónica} con la ejecución.
    \item \textbf{Trap explícita / syscall}: una instrucción
        especial (\texttt{syscall} en x86, \texttt{ecall} en
        RISC-V, \texttt{svc} en ARM) que el proceso ejecuta para
        pedir al kernel un servicio.
\end{itemize}

Las tres usan el \emph{mismo} mecanismo de hardware: la CPU detecta
el evento, salva el estado mínimo y salta a un punto fijo del kernel.

\subsection*{Flujo completo paso a paso}

\paragraph{Fase 1 --- Hardware atiende el trap.}

\begin{enumerate}[leftmargin=*]
    \item Si las interrupciones están habilitadas, la CPU acepta el
        evento.
    \item \textbf{Salva en registros especiales} el estado mínimo:
        el \texttt{PC} actual (para saber a dónde volver), el modo
        de ejecución previo (user/kernel), y la causa del trap
        (\texttt{scause} en RISC-V, vector number en x86).
    \item Cambia a \textbf{modo supervisor} si no estaba ya.
    \item Salta a una dirección fija del kernel definida en el
        \emph{trap vector} (\texttt{stvec} en RISC-V, IDT en x86,
        \texttt{VBAR} en ARM).
\end{enumerate}

\paragraph{Fase 2 --- Trampolín y guardado completo del contexto.}

El primer código que se ejecuta del kernel es un pequeño
\emph{trampolín} en assembly (\texttt{uservec} en xv6). Su tarea es:

\begin{enumerate}[leftmargin=*]
    \item Conmutar a un stack del kernel (no del proceso).
    \item Salvar \textbf{todos los registros} del CPU (de propósito
        general, de estado, etc.) en una estructura llamada
        \textbf{\texttt{trapframe}} en el stack del kernel.
    \item Cargar la tabla de páginas del kernel (si era una trap
        desde modo usuario).
    \item Saltar a una función C principal (e.g., \texttt{usertrap}).
\end{enumerate}

\paragraph{Fase 3 --- Despacho según la causa.}

La función C principal lee el campo de causa (e.g., \texttt{scause})
y despacha al \emph{handler} correspondiente:

\begin{itemize}[leftmargin=*,nosep]
    \item Syscall $\to$ \texttt{syscall()}, que extrae el número
        de syscall del trapframe y llama a la función
        correspondiente (\texttt{sys\_read}, \texttt{sys\_write},
        \texttt{sys\_fork}, etc.).
    \item Page fault $\to$ \texttt{page\_fault\_handler()}.
    \item Timer / IRQ de disco $\to$ ISR del driver correspondiente.
    \item Excepción no recuperable (división por cero, etc.) $\to$
        enviar \texttt{SIGSEGV}/\texttt{SIGFPE} al proceso.
\end{itemize}

\paragraph{Fase 4 --- Posible cambio de proceso.}

Después del manejo, el kernel puede invocar al scheduler:

\begin{itemize}[leftmargin=*,nosep]
    \item Si fue interrupción del timer y el proceso usó su quantum
        $\to$ \texttt{yield()}: el scheduler elige otro proceso.
        Cambio de contexto.
    \item Si fue un syscall bloqueante (\texttt{read} sobre disco):
        el proceso pasa a \texttt{SLEEPING} y el scheduler elige
        otro.
\end{itemize}

Eventualmente, el scheduler vuelve a planificar el proceso original
(en algún momento futuro).

\paragraph{Fase 5 --- Retorno.}

Cuando es momento de regresar al proceso interrumpido:

\begin{enumerate}[leftmargin=*]
    \item Se carga la tabla de páginas del proceso.
    \item Se ejecuta el trampolín de retorno (\texttt{userret} en
        xv6), que restaura todos los registros desde el trapframe.
    \item Se ejecuta la instrucción atómica de retorno
        (\texttt{sret} en RISC-V, \texttt{iret} en x86,
        \texttt{eret} en ARM). Esta instrucción restaura el
        \texttt{PC} desde \texttt{sepc}, el modo a usuario y
        rehabilita las interrupciones, todo en un solo paso.
\end{enumerate}

El proceso reanuda como si nada hubiera pasado.

\subsection*{El rol del \emph{trapframe}}

El \texttt{trapframe} es una \textbf{estructura de datos en el stack
del kernel} del proceso, asociada al PCB, que conserva el estado
completo del proceso al momento de la trap.

\textbf{Estructura típica} (xv6/RISC-V):

\begin{lstlisting}[style=cstyle]
struct trapframe {
    uint64 kernel_satp;   // dir.  tabla de paginas del kernel
    uint64 kernel_sp;     // stack pointer del kernel
    uint64 kernel_trap;   // direccion de usertrap()
    uint64 epc;           // PC en el momento del trap (sepc)
    uint64 kernel_hartid; // CPU actual

    /* Registros de proposito general */
    uint64 ra, sp, gp, tp;
    uint64 t0, t1, t2, t3, t4, t5, t6;
    uint64 s0, s1, ..., s11;
    uint64 a0, a1, ..., a7;
};
\end{lstlisting}

\textbf{Funciones del trapframe:}

\begin{enumerate}[leftmargin=*]
    \item \textbf{Salvar el estado del proceso}: cuando ocurre un
        trap, el trampolín copia todos los registros al trapframe.
        Esto permite que el handler del kernel use los registros
        libremente sin perder el estado del proceso.

    \item \textbf{Pasar argumentos al handler}: en una syscall, los
        argumentos del usuario están en registros (\texttt{a0--a7}
        en RISC-V). El handler los lee del trapframe.

    \item \textbf{Devolver resultados al proceso}: para que una
        syscall retorne un valor, el handler escribe el resultado en
        el campo correspondiente del trapframe (e.g., \texttt{a0}).
        Cuando el trampolín de retorno restaura los registros, el
        proceso ``ve'' el resultado en su registro \texttt{a0}.

    \item \textbf{Soportar cambios de contexto}: si el handler decide
        cambiar de proceso, el trapframe del proceso saliente queda
        guardado; al volver a planificarlo, se restaura desde su
        propio trapframe.

    \item \textbf{Auditoría / debugging}: el contenido del trapframe
        contiene exactamente la \emph{imagen registral} del proceso
        en el momento de la trap, útil para diagnóstico (e.g.,
        para reportar el \texttt{PC} donde ocurrió un page fault).
\end{enumerate}

\subsection*{Diagrama del flujo}

\begin{center}
\begin{tikzpicture}[
    every node/.style={font=\ttfamily\small,align=center},
    box/.style={draw,rounded corners,minimum width=3.2cm,minimum height=0.85cm},
    user/.style={box,fill=c1!25},
    hw/.style={box,fill=c2!25},
    kern/.style={box,fill=c3!25},
    arr/.style={->,thick}
]
\node[user] (P1) at (0,3) {Proceso \\ ejecutando};
\node[hw]   (HW) at (0,1.5) {Hardware salva: \\ \scriptsize PC, modo, causa};
\node[kern] (TR) at (0,0)   {Trampolín: \\ \scriptsize guarda regs en trapframe};
\node[kern] (HD) at (5,0)   {usertrap(): \\ \scriptsize despacha handler};
\node[kern] (SC) at (5,1.5) {sys\_X() o \\ page\_fault\_handler() o ISR};
\node[kern] (RT) at (5,3)   {trampolín retorno: \\ \scriptsize restaura regs};
\node[user] (P2) at (0,4.5) {Proceso \\ reanuda};

\draw[arr] (P1) -- node[right,font=\scriptsize]{trap} (HW);
\draw[arr] (HW) -- (TR);
\draw[arr] (TR) -- (HD);
\draw[arr] (HD) -- (SC);
\draw[arr] (SC) -- (RT);
\draw[arr] (RT) -| node[right,font=\scriptsize,xshift=-2cm]{sret/iret} (P2);
\end{tikzpicture}
\end{center}

\begin{respuesta}
\textbf{Pasos generales:} (1) la CPU detecta el trap, salva un
estado mínimo (PC, modo, causa) y salta al \emph{trap vector}; (2)
un trampolín en assembly del kernel guarda todos los registros del
proceso en el \textbf{trapframe} (estructura asociada al PCB) y
cambia al stack del kernel; (3) una función C en el kernel despacha
al \emph{handler} adecuado según la causa; (4) el handler procesa
(eventualmente cambia de contexto); (5) al volver, el trampolín de
retorno restaura los registros desde el trapframe y ejecuta la
instrucción atómica de retorno (\texttt{sret}/\texttt{iret}/\texttt{eret})
que devuelve el control al proceso en modo usuario.

\textbf{Rol del trapframe:} es el ``snapshot'' completo del estado
del proceso al ocurrir el trap. Sirve para (a) que el handler del
kernel pueda usar los registros sin perder el contexto del proceso,
(b) leer los argumentos de una syscall, (c) escribir el resultado
de la syscall (que será visible al proceso cuando reanude), y (d)
soportar cambios de contexto: cada proceso tiene su propio
trapframe asociado a su PCB.
\end{respuesta}

% =====================================================================
\section{Ejercicio 2 --- Lock reentrante sin \emph{selflock} (2 ptos)}
% =====================================================================

\textbf{Enunciado.} Un proceso o el kernel produce un \emph{selflock}
si invoca a \texttt{lock(\&lk)} dos veces sin que en el medio haya un
\texttt{unlock(\&lk)}. Proponga una implementación en pseudo-código
de \texttt{lock}/\texttt{unlock} que permita invocar repetidamente a
\texttt{lock(\&lk)} sin que se produzca un selflock. Asuma que
existen \texttt{cpuid()} y \texttt{pid()} que retornan el ID de la
CPU y el PID actual.

\subsection*{Análisis: lock reentrante (recursive lock)}

La idea es identificar al \emph{dueño actual} del lock por la
combinación \texttt{(cpuid, pid)} y mantener un \textbf{contador}:

\begin{itemize}[leftmargin=*,nosep]
    \item Si el lock está libre, lo tomamos normalmente y guardamos
        nuestra identidad.
    \item Si el lock está tomado y \emph{nosotros somos el dueño},
        simplemente incrementamos el contador y retornamos sin
        bloquearnos.
    \item Si el lock está tomado por \emph{otro}, esperamos
        (\emph{spin}) como un lock normal.
    \item \texttt{unlock} decrementa el contador. Solo cuando llega
        a $0$ libera realmente el lock.
\end{itemize}

Esto convierte el lock en \textbf{reentrante}: el mismo dueño puede
tomarlo varias veces de forma anidada sin selflock.

\subsection*{Pseudo-código}

\begin{lstlisting}[style=cstyle]
typedef struct {
    int        locked;       /* 0 = libre, 1 = tomado */
    int        owner_cpu;    /* CPU del dueno actual */
    int        owner_pid;    /* PID del dueno actual */
    int        depth;        /* veces que el mismo dueno lo tomo */
} lock_t;

void lock(lock_t *lk) {
    int my_cpu = cpuid();
    int my_pid = pid();

    /* 1) Si YA somos los duenos, simplemente incrementamos. */
    if (lk->locked == 1 &&
        lk->owner_cpu == my_cpu &&
        lk->owner_pid == my_pid) {
        lk->depth++;
        return;
    }

    /* 2) No lo tenemos: intentamos adquirirlo atomicamente. */
    while (atomic_swap(&lk->locked, 1) == 1) {
        /* Si lo tenia otro dueno, seguimos esperando.
         * Si lo tenia el mismo dueno, no llegamos a este while
         * (volvimos en el if anterior). */
        ;     /* spin-wait */
    }

    /* 3) Lo conseguimos: marcamos nuestra identidad. */
    lk->owner_cpu = my_cpu;
    lk->owner_pid = my_pid;
    lk->depth     = 1;
}

void unlock(lock_t *lk) {
    int my_cpu = cpuid();
    int my_pid = pid();

    /* Verificacion de coherencia. */
    if (lk->locked != 1 ||
        lk->owner_cpu != my_cpu ||
        lk->owner_pid != my_pid) {
        panic("unlock por un thread que no es el dueno");
    }

    lk->depth--;
    if (lk->depth == 0) {
        lk->owner_cpu = -1;
        lk->owner_pid = -1;
        lk->locked    = 0;    /* liberacion atomica */
    }
    /* Si depth > 0, todavia hay locks anidados del mismo dueno;
     * el lock NO se libera. */
}
\end{lstlisting}

\subsection*{Por qué funciona}

\begin{itemize}[leftmargin=*]
    \item \textbf{Sin selflock}: si el mismo \texttt{(cpu, pid)}
        llama dos veces a \texttt{lock}, la segunda invocación cae
        en el primer \texttt{if} y simplemente incrementa
        \texttt{depth}. No espera en el \texttt{while}.
    \item \textbf{Exclusión mutua con otros}: si otro
        \texttt{(cpu, pid)} intenta tomar, ve \texttt{locked = 1}
        con dueño distinto, no entra al \texttt{if} y se va al
        \texttt{while}, donde espera.
    \item \textbf{Liberación correcta}: cada \texttt{lock} hace
        \texttt{depth++}, cada \texttt{unlock} hace
        \texttt{depth--}. Solo cuando \texttt{depth} llega a $0$ se
        libera \texttt{locked} de verdad.
\end{itemize}

\subsection*{Ejemplo de uso}

\begin{lstlisting}[style=cstyle]
void funcion_externa() {
    lock(&lk);          /* depth pasa a 1 */
    funcion_interna();
    unlock(&lk);        /* depth pasa a 0, libera */
}

void funcion_interna() {
    lock(&lk);          /* SAME owner; depth pasa a 2; no se bloquea */
    /* ... */
    unlock(&lk);        /* depth pasa a 1 */
}
\end{lstlisting}

Sin lock reentrante, \texttt{lock(\&lk)} dentro de
\texttt{funcion\_interna} se bloquearía contra sí misma
(\emph{selflock}). Con el lock reentrante propuesto, funciona
correctamente.

\begin{respuesta}
La clave es identificar al \emph{dueño actual} con \texttt{(cpuid,
pid)} y mantener un contador \texttt{depth}. \texttt{lock} verifica
primero si el llamante ya es dueño; si lo es, incrementa
\texttt{depth} y retorna. \texttt{unlock} decrementa y solo libera
realmente cuando \texttt{depth} llega a $0$.

\textbf{Importante}: la operación de adquisición debe ser atómica
(\texttt{atomic\_swap}, \texttt{cmpxchg}, LL/SC) para garantizar
exclusión mutua con otros dueños en un multiprocesador.
\end{respuesta}

% =====================================================================
\section{Ejercicio 3 --- FIFO vs Round Robin (2 ptos)}
% =====================================================================

\textbf{Enunciado.} Describa la diferencia entre el algoritmo de
planificación FIFO y Round Robin.

\subsection*{FIFO (First In First Out, también FCFS)}

\textbf{Idea.} Una sola cola \texttt{READY}. Los procesos se
encolan al llegar; el primero en arribar es el primero en ejecutar.
Una vez asignada la CPU, el proceso \textbf{ejecuta hasta terminar
o bloquearse} (no preemption).

\textbf{Propiedades:}

\begin{itemize}[leftmargin=*,nosep]
    \item \textbf{No-preemptivo}: ningún proceso es desalojado por
        agotar tiempo.
    \item Implementación: una cola simple FIFO.
    \item Métricas: turnaround promedio depende fuertemente del
        orden de llegada y de los tamaños de las ráfagas.
\end{itemize}

\textbf{Convoy effect.} Si un proceso CPU-bound largo llega primero,
todos los que vienen detrás esperan hasta que termine, aunque sean
cortos. El tiempo de espera promedio crece mucho.

\textbf{Diagrama de Gantt típico} (P1: 10, P2: 1, P3: 1, en orden de
llegada):

\begin{center}
\begin{tikzpicture}[xscale=0.5, yscale=0.7, every node/.style={font=\ttfamily\small}]
\draw[fill=c1!50] (0,0)  rectangle (10,1) node[midway,white,font=\bfseries]{P1};
\draw[fill=c2!50] (10,0) rectangle (11,1) node[midway,white,font=\bfseries\scriptsize]{P2};
\draw[fill=c3!50] (11,0) rectangle (12,1) node[midway,white,font=\bfseries\scriptsize]{P3};
\foreach \x in {0,10,11,12}
    \draw (\x,0) -- (\x,-0.2) node[below,font=\scriptsize]{\x};
\end{tikzpicture}
\end{center}

\noindent P2 y P3 esperan $\sim 10$ unidades aunque solo necesitaran
$1$. Turnaround promedio: $(10 + 11 + 12)/3 \approx 11$.

\subsection*{Round Robin}

\textbf{Idea.} Una cola \texttt{READY}, pero a cada proceso se le
asigna un \textbf{quantum} (\emph{time slice}, $q$). Si al expirar el
quantum el proceso aún no terminó, se lo \textbf{desaloja} (preempt)
y se reencola al final. El scheduler toma el siguiente.

\textbf{Propiedades:}

\begin{itemize}[leftmargin=*,nosep]
    \item \textbf{Preemptivo}: el timer fuerza el desalojo al
        expirar el quantum.
    \item Implementación: cola FIFO + temporizador.
    \item Métricas: garantiza un tiempo de respuesta acotado por
        $N \times q$ (con $N$ procesos en \texttt{READY}).
\end{itemize}

\textbf{Sin convoy effect.} Aunque haya un proceso largo, los demás
obtienen CPU cada $\leq N \cdot q$ unidades. Bueno para sistemas
interactivos.

\textbf{Diagrama de Gantt típico} (mismos procesos, $q = 2$):

\begin{center}
\begin{tikzpicture}[xscale=0.5, yscale=0.7, every node/.style={font=\ttfamily\small}]
\draw[fill=c1!50] (0,0)  rectangle (2,1) node[midway,white,font=\bfseries\scriptsize]{P1};
\draw[fill=c2!50] (2,0)  rectangle (3,1) node[midway,white,font=\bfseries\scriptsize]{P2};
\draw[fill=c3!50] (3,0)  rectangle (4,1) node[midway,white,font=\bfseries\scriptsize]{P3};
\draw[fill=c1!50] (4,0)  rectangle (6,1) node[midway,white,font=\bfseries\scriptsize]{P1};
\draw[fill=c1!50] (6,0)  rectangle (8,1) node[midway,white,font=\bfseries\scriptsize]{P1};
\draw[fill=c1!50] (8,0)  rectangle (10,1) node[midway,white,font=\bfseries\scriptsize]{P1};
\draw[fill=c1!50] (10,0) rectangle (12,1) node[midway,white,font=\bfseries\scriptsize]{P1};
\foreach \x in {0,2,3,4,6,8,10,12}
    \draw (\x,0) -- (\x,-0.2) node[below,font=\scriptsize]{\x};
\end{tikzpicture}
\end{center}

\noindent P2 termina en $t = 3$ y P3 en $t = 4$, muchísimo mejor que
con FIFO. P1 termina al final, pero no hace esperar a los demás.

\subsection*{Comparación}

\begin{center}
\small
\begin{tabular}{@{}p{4.5cm}p{4.5cm}p{4.5cm}@{}}
\toprule
Característica & FIFO (FCFS) & Round Robin \\
\midrule
Preemptivo                    & No                    & Sí (al expirar quantum) \\
Estructura                    & cola FIFO             & cola FIFO $+$ timer \\
Tiempo de respuesta           & Malo (varía mucho)    & Acotado por $N \cdot q$ \\
Convoy effect                 & Sufre                 & Lo evita \\
Cambios de contexto           & Pocos                 & Muchos (cada quantum) \\
Apto para interactivos        & No                    & Sí \\
Apto para batch               & Sí                    & Solo si $q$ es grande \\
Caso $q \to \infty$           & ---                   & Degenera en FIFO \\
Caso $q \to 0$                & ---                   & Overhead domina \\
\bottomrule
\end{tabular}
\end{center}

\begin{respuesta}
\textbf{Diferencia esencial}: FIFO es \textbf{no preemptivo} y
ejecuta los procesos en estricto orden de llegada hasta que
terminen. Round Robin es \textbf{preemptivo} y le asigna a cada
proceso un cuanto de tiempo (\emph{quantum}); al expirar el
quantum, el proceso vuelve al final de la cola.

\textbf{Consecuencia práctica}: RR evita el \emph{convoy effect} y
garantiza tiempo de respuesta acotado, ideal para sistemas
interactivos. FIFO es más simple y tiene menos cambios de contexto,
ideal para sistemas \emph{batch} con cargas predecibles. El
\emph{quantum} de RR es el parámetro que regula el compromiso
respuesta vs.\ overhead: si $q \to \infty$, RR se vuelve FIFO; si
$q \to 0$, el overhead de cambios de contexto domina.
\end{respuesta}

% =====================================================================
\section{Ejercicio 4 --- V/F sobre memoria virtual (2 ptos)}
% =====================================================================

\subsection*{4.a --- ``El paginado define espacios lógicos a cada proceso; la segmentación no''}

\begin{respuesta}
\textbf{F}. \emph{Ambas} técnicas, paginado y segmentación,
permiten definir un espacio lógico de direcciones para cada proceso.
La diferencia es la \emph{unidad} de mapeo:

\begin{itemize}[leftmargin=*,nosep]
    \item \textbf{Paginado}: el espacio se divide en páginas de
        tamaño fijo. La tabla de páginas mapea \texttt{vpn} a
        \texttt{ppn}.
    \item \textbf{Segmentación}: el espacio se divide en segmentos
        de tamaño variable (código, datos, stack, etc.) con base y
        límite. La tabla de descriptores (GDT/LDT en x86) mapea
        cada segmento a su ubicación física.
\end{itemize}

Lo correcto sería que \emph{ambas} ofrecen aislamiento entre
procesos. De hecho, x86 históricamente usó segmentación
\emph{antes} de tener paginado. La segmentación pura cayó en
desuso (es menos eficiente para memoria virtual; ver ítem b), pero
\emph{sí} provee espacios lógicos.
\end{respuesta}

\subsection*{4.b --- ``El paginado es más adecuado que la segmentación para memoria virtual''}

\begin{respuesta}
\textbf{V}. La paginación tiene varias ventajas críticas para
memoria virtual:

\begin{itemize}[leftmargin=*,nosep]
    \item \textbf{Tamaño uniforme}: todas las páginas miden lo
        mismo. El swap-in/swap-out es de tamaño fijo, simple y
        predecible.
    \item \textbf{Sin fragmentación externa}: cualquier frame
        físico libre sirve. En segmentación, segmentos de tamaño
        variable producen huecos no aprovechables.
    \item \textbf{Granularidad fina}: con \emph{demand paging}, se
        carga solo lo que se accede. Con segmentación, habría que
        traer un segmento entero (puede ser MB).
    \item \textbf{TLB y MMU simples}: el hardware de paginado es
        más sencillo y rápido que el de segmentación + paginado
        combinado.
\end{itemize}

Por eso todos los SO modernos usan \textbf{paginado} (o paginado
combinado con segmentación residual, como x86), y los segmentos
``puros'' son cosa del pasado.
\end{respuesta}

\subsection*{4.c --- ``COW puede usarse para implementar \texttt{fork()} eficientemente''}

\begin{respuesta}
\textbf{V}. Es exactamente la motivación de \emph{copy-on-write}.
La implementación de \texttt{fork()} con COW es:

\begin{enumerate}[leftmargin=*,nosep]
    \item Al hacer \texttt{fork()}, el kernel \emph{no copia} las
        páginas del padre. En cambio, crea para el hijo una nueva
        tabla de páginas que apunta a los mismos \emph{frames
        físicos} del padre.
    \item Las páginas se marcan como \textbf{solo lectura} en
        ambas tablas (padre e hijo).
    \item Mientras ambos solo lean, comparten físicamente (gratis).
    \item Cuando uno \emph{escribe} a una página compartida, el
        hardware levanta un \emph{page fault}. El kernel atiende:
        copia la página afectada a un nuevo frame, actualiza la
        tabla de quien escribió para que apunte ahí (writable), y
        deja al otro con la copia original.
\end{enumerate}

\textbf{Por qué es eficiente}: en el caso típico (\texttt{fork +
exec}), el hijo no escribe nada del padre antes de
\texttt{exec()}, así que NUNCA se copia ninguna página. El costo de
\texttt{fork} se reduce a copiar la tabla de páginas y unos
metadatos, no varios MB de RAM.
\end{respuesta}

\subsection*{4.d --- ``\texttt{mmap()} puede usarse para implementar \texttt{exec()} eficientemente''}

\begin{respuesta}
\textbf{V}. Esto es exactamente lo que hacen los SO modernos. La
implementación de \texttt{exec()} con \texttt{mmap}:

\begin{enumerate}[leftmargin=*,nosep]
    \item En vez de leer todo el ejecutable a memoria, \texttt{exec}
        hace \texttt{mmap} del archivo ejecutable en el espacio de
        direcciones del proceso.
    \item Las páginas se marcan como ``presentes en disco'' pero
        no en RAM (\texttt{V = 0} en la \texttt{pte}, con
        información en otro campo de dónde leer).
    \item Cuando el proceso empieza a ejecutar instrucciones,
        cualquier acceso a una página todavía no cargada produce un
        \emph{page fault}.
    \item El kernel atiende el fault leyendo \emph{esa} página del
        disco.
\end{enumerate}

\textbf{Por qué es eficiente}:

\begin{itemize}[leftmargin=*,nosep]
    \item Solo se lee del disco lo que efectivamente se ejecuta o
        accede. Muchas funciones del binario nunca se llaman, y sus
        páginas nunca se cargan.
    \item Si varias instancias del mismo binario corren en
        paralelo, las páginas de \texttt{.text} (read-only) se
        comparten físicamente.
    \item El kernel administra una \emph{page cache} unificada;
        las páginas leídas pueden quedar en RAM aunque otro proceso
        cierre el binario, facilitando que el siguiente
        \texttt{exec()} del mismo programa sea instantáneo.
\end{itemize}
\end{respuesta}

% =====================================================================
\section{Ejercicio 5 --- Paginación: tabla y page faults (2 ptos)}
% =====================================================================

\textbf{Enunciado.} En una arquitectura de 32 bits, espacio de
direcciones de 4\,GB y páginas de 4\,KB, un proceso mapeado desde la
dirección 0 tiene:

\begin{itemize}[leftmargin=*,nosep]
    \item $7$\,KB de código.
    \item $3.5$\,KB de datos globales.
    \item $2$ páginas de stack a partir de \texttt{0x7fe000}.
\end{itemize}

\textbf{(a)} Dar el esquema de la tabla de páginas del proceso.
\textbf{(b)} ¿La dirección \texttt{0x381A} produce page fault?
\textbf{(c)} Ídem para \texttt{0x7ff100}.

\subsection*{Análisis de parámetros}

\begin{itemize}[leftmargin=*,nosep]
    \item Página: $4$\,KB $= 2^{12} = $ \texttt{0x1000} bytes.
    \item Espacio total: $4$\,GB $= 2^{32}$ bytes.
    \item Total de páginas posibles: $2^{32} \div 2^{12} = 2^{20} =
        1\,048\,576$ páginas. La tabla tiene $2^{20}$ entradas
        (típicamente representadas como un árbol multinivel para no
        ocupar $4$\,MB por proceso).
    \item Direcciones virtuales: $20$ bits altos = vpn, $12$ bits
        bajos = offset.
\end{itemize}

\subsection*{Layout del proceso}

\textbf{Páginas por área (con padding a frontera de página, como hace
el loader ELF):}

\begin{itemize}[leftmargin=*,nosep]
    \item \textbf{Código} (7\,KB): $\lceil 7 / 4 \rceil = 2$ páginas
        $\to$ vpn $0$ y $1$ (rango virtual
        \texttt{[0x0000, 0x2000)}).
    \item \textbf{Datos globales} (3.5\,KB): $\lceil 3.5 / 4 \rceil
        = 1$ página $\to$ vpn $2$ (rango \texttt{[0x2000, 0x3000)}).
    \item \textbf{Stack}: $2$ páginas a partir de
        \texttt{0x7fe000}. Como cada página es $0x1000$, las
        páginas son
        \texttt{vpn = 0x7fe000 / 0x1000 = 0x7fe} $= 2046$ y
        \texttt{vpn = 0x7ff} $= 2047$.
\end{itemize}

\subsection*{5.a --- Esquema de la tabla de páginas}

\begin{center}
\footnotesize
\begin{tabular}{@{}llccp{2.5cm}p{3cm}@{}}
\toprule
vpn (hex)    & rango virtual               & V & R/W/X & Área              & Notas \\
\midrule
\texttt{0x000} & \texttt{[0x0000, 0x1000)} & 1 & R + X & código pág.\ 1    & 4\,KB \\
\texttt{0x001} & \texttt{[0x1000, 0x2000)} & 1 & R + X & código pág.\ 2    & 3\,KB + 1\,KB pad \\
\texttt{0x002} & \texttt{[0x2000, 0x3000)} & 1 & R + W & datos globales    & 3.5\,KB + 0.5\,KB pad \\
\texttt{0x003}--\texttt{0x7fd} & \texttt{[0x3000, 0x7fe000)} & 0 & --- & gap inválido & 2043 páginas \\
\texttt{0x7fe} & \texttt{[0x7fe000, 0x7ff000)} & 1 & R + W & stack pág.\ 1     & \\
\texttt{0x7ff} & \texttt{[0x7ff000, 0x800000)} & 1 & R + W & stack pág.\ 2     & \\
$\geq$\texttt{0x800} & \texttt{[0x800000, \ldots)} & 0 & --- & inválidas         & resto \\
\bottomrule
\end{tabular}
\end{center}

Solo \textbf{5 entradas válidas} (vpn $0, 1, 2, 0x7fe, 0x7ff$) de
$2^{20}$ posibles. El resto está inválido (V = 0). Esto justifica
por qué la tabla se implementa como un árbol multinivel: solo los
nodos correspondientes a regiones efectivamente usadas se reservan
en RAM.

\subsection*{Diagrama del espacio de direcciones}

\begin{center}
\begin{tikzpicture}[every node/.style={font=\ttfamily\small},
    sec/.style={draw,minimum width=5cm,minimum height=0.6cm,align=center}]
\node[sec,fill=c1!25] (code1) at (0,4.5) {vpn 0: código};
\node[sec,fill=c1!25,below=0pt of code1] (code2) {vpn 1: código (+pad)};
\node[sec,fill=c2!25,below=0pt of code2] (data) {vpn 2: datos (+pad)};
\node[sec,fill=cred!15,below=0.4cm of data,minimum height=1.2cm] (gap) {GAP \\ vpn 3 a 2045 \\ (\texttt{[0x3000, 0x7fe000)})};
\node[sec,fill=c3!25,below=0.4cm of gap] (stk1) {vpn 2046: stack};
\node[sec,fill=c3!25,below=0pt of stk1] (stk2) {vpn 2047: stack};

\node[anchor=west,xshift=0.4cm] at (code1.east)  {\texttt{0x00000000}};
\node[anchor=west,xshift=0.4cm] at (code2.east)  {\texttt{0x00001000}};
\node[anchor=west,xshift=0.4cm] at (data.east)   {\texttt{0x00002000}};
\node[anchor=west,xshift=0.4cm] at (data.south east) [yshift=-2pt] {\texttt{0x00003000}};
\node[anchor=west,xshift=0.4cm] at (stk1.east)   {\texttt{0x007fe000}};
\node[anchor=west,xshift=0.4cm] at (stk2.east)   {\texttt{0x007ff000}};
\node[anchor=west,xshift=0.4cm] at (stk2.south east) [yshift=-2pt] {\texttt{0x00800000}};
\end{tikzpicture}
\end{center}

\subsection*{5.b --- ¿\texttt{0x381A} produce page fault?}

\textbf{Descomposición} (20 bits vpn, 12 bits offset):

\quad \texttt{0x381A} en 32 bits = \texttt{0x0000381A}

\quad \texttt{0x381A} $\div$ \texttt{0x1000} $= 3$ (cociente)

\quad \texttt{0x381A} $\bmod$ \texttt{0x1000} $= $ \texttt{0x81A}

Por tanto: \texttt{vpn = 3}, \texttt{offset = 0x81A}.

\textbf{Consulta a la tabla}: la entrada \texttt{vpn = 3} cae en el
gap inválido (entre vpn 3 y vpn 2045). \texttt{V = 0}.

\begin{respuesta}
\textbf{(5.b) SÍ produce page fault.} La dirección \texttt{0x381A}
está en \texttt{vpn = 3}, primera página después del área de datos.
Esa página no está mapeada (V = 0). El hardware genera un page
fault y el kernel envía \texttt{SIGSEGV} al proceso.
\end{respuesta}

\subsection*{5.c --- ¿\texttt{0x7ff100} produce page fault?}

\textbf{Descomposición}:

\quad \texttt{0x7ff100} $\div$ \texttt{0x1000} $= $ \texttt{0x7ff} $= 2047$

\quad \texttt{0x7ff100} $\bmod$ \texttt{0x1000} $= $ \texttt{0x100}

Por tanto: \texttt{vpn = 0x7ff} $= 2047$, \texttt{offset = 0x100}.

\textbf{Consulta a la tabla}: la entrada \texttt{vpn = 0x7ff}
corresponde a la \textbf{segunda página de stack}. Está válida
(\texttt{V = 1}, \texttt{R + W}).

\begin{respuesta}
\textbf{(5.c) NO produce page fault.} La dirección \texttt{0x7ff100}
está en \texttt{vpn = 0x7ff = 2047}, que es la \emph{segunda página
del stack} del proceso (la que arranca en \texttt{0x7ff000}). Está
mapeada, con permisos R+W. El acceso es legítimo.
\end{respuesta}

\subsection*{Verificación numérica}

\begin{itemize}[leftmargin=*,nosep]
    \item Stack comienza en \texttt{0x7fe000} (\texttt{vpn 0x7fe}).
        $0\text{x}7\text{fe} \cdot 0\text{x}1000 = 0\text{x}7\text{fe}000$. ✓
    \item Stack ocupa $2$ páginas: \texttt{vpn 0x7fe} y
        \texttt{vpn 0x7ff}, rango total
        \texttt{[0x7fe000, 0x800000)}.
    \item \texttt{0x7ff100} cae claramente dentro de ese rango. ✓
\end{itemize}

% =====================================================================
\section*{Resumen del recuperatorio}
\addcontentsline{toc}{section}{Resumen del recuperatorio}
% =====================================================================

\begin{center}
\small
\begin{tabular}{@{}clp{9cm}@{}}
\toprule
\# & Tema & Respuesta sintética \\
\midrule
1 & Traps + trapframe & Hardware salva estado mínimo, trampolín guarda regs en trapframe, dispatch a handler C, retorno con \texttt{sret}. Trapframe = snapshot completo del estado del proceso \\
2 & Lock reentrante  & Guardar \texttt{(cpu, pid)} del dueño + contador \texttt{depth}; si soy el dueño, incremento sin bloquear; \texttt{unlock} decrementa y libera al llegar a $0$ \\
3 & FIFO vs RR       & FIFO: no preemptivo, sufre convoy; RR: preemptivo con quantum, evita convoy, mejor para interactivos a costa de overhead de cambios de contexto \\
4 & V/F memoria      & a:F (segmentación también define espacios), b:V (uniformidad, no frag.\ externa), c:V (COW), d:V (mmap del ejecutable con demand paging) \\
5 & Paginación num.  & (a) 5 entradas válidas: vpn 0,1,2 + 2046, 2047; (b) \texttt{0x381A} vpn=3 (gap) → page fault; (c) \texttt{0x7ff100} vpn=2047 (stack) → NO fault \\
\bottomrule
\end{tabular}
\end{center}

\end{document}
